\begin{statement}[C.7.11; Proved: the complexified route and its exact obstructions]
Use the curve, point and real logarithm of C.7.9:
\[
 E:y^2+y=x^3-38x+90,\qquad Q=(4,-2)=\exp_Eu,
 \qquad 0<u<\omega_1/2.
\]
Put $\beta=i\sqrt{19}$, $k=\pi/A$, $s_2=-2$, and
\[
 \theta(z)=\sigma(z)e^{-s_2z^2/2},\qquad
 D(z)=\zeta(z)-s_2z,\qquad
 \mathcal F(z,w)=\sqrt{19}\,\theta(z)\theta(w)e^{-kzw},
 \qquad H=\mathcal F^2.
\]
Then $\mathcal F(z,\bar z)=N(z)^2$, and $\mathcal F$ is an
entire function of two complex variables of order at most two.
It is periodic under $(z,w)\mapsto(z+\omega,w+\bar\omega)$,
$\omega\in\Lambda$. The subgroup
\[
 \Gamma=\{(\alpha u+\omega,\bar\alpha u+\bar\omega):
                   \alpha\in\Z[\beta],\ \omega\in\Lambda\}
 \subset\C^2
\]
has rank four over $\Z$ and lies in the real plane
$\{(z,\bar z):z\in\C\}$. If $N(Q)$ is algebraic, then the
values of $H$ on $\Gamma$ lie in one number field.

For every nonzero $\alpha\in\Z[\beta]$ there is an algebraic
number $a_\alpha$ such that
\[
 D(\alpha u)=\bar\alpha D(u)+a_\alpha.
\]
At $p_\alpha=(\alpha u,\bar\alpha u)$ the algebraic constant
direction $V_\alpha=\alpha\partial_z-\bar\alpha\partial_w$
satisfies
\begin{align*}
 V_\alpha\log\mathcal F(p_\alpha)
    &=\alpha a_\alpha-\bar\alpha\bar a_\alpha
                         \in\overline{\Q},\\
 V_\alpha^2\log\mathcal F(p_\alpha)
    &=-\alpha^2\{\wp(\alpha u)+s_2\}
      -\bar\alpha^2\{\wp(\bar\alpha u)+s_2\}
      +2\alpha\bar\alpha k.
\end{align*}
The number $k$ is transcendental. Consequently the second
directional derivative of $H$ at $p_\alpha$ is transcendental
whenever $H(p_\alpha)$ is algebraic. In particular,
\[
 (\partial_z-\partial_w)H(u,u)=0,\qquad
 \frac{(\partial_z-\partial_w)^2H(u,u)}{H(u,u)}=4k-8.
\]
Thus the cancelling first derivative does not give algebraic jets
of all orders. Moreover the rank-four set does not satisfy the
explicit derivative-free finite-order criterion audited below.
These are obstructions to the stated applications of those
criteria; they do not decide the algebraicity of $N(Q)$.
\end{statement}

\begin{proof}
The real model and its algebraic differential give real sigma
coefficients and a conjugation-stable lattice. In particular
$\overline{\theta(z)}=\theta(\bar z)$. The discriminant is
$-19^3$, so C.7.1 gives the asserted diagonal identity.
The sigma function and the exponential of the quadratic
polynomial $-kzw$ have order at most two. From C.7.6,
\[
 \theta(z+\omega)=\epsilon(\omega)
         e^{k\bar\omega(z+\omega/2)}\theta(z),
 \qquad \epsilon(\bar\omega)=\epsilon(\omega)\in\{1,-1\}.
\]
Multiplying this formula and its conjugate version, the two
signs multiply to one and the exponent cancels the change of
$-kzw$. This proves periodicity on all of $\C^2$, not only
on the real diagonal.

If $\alpha u+\omega=0$, then $[\alpha]Q=O$.
Nontorsion of $Q$ and the fact that a nonzero elliptic
endomorphism is an isogeny imply $\alpha=0$, and then
$\omega=0$. Thus the rank-two module $\Z[\beta]u$ and the
rank-two period lattice are independent over $\Z$.
This proves the rank assertion, with no extra algebraic
independence assumption on $u$ and the periods.

For $d=\alpha\bar\alpha$, the squared CM law C.7.6 gives
the exact fixed-field identity
\begin{equation}\label{eq:cc-fixedfield}
 H(\alpha u,\bar\alpha u)
   =19^{1-d}F_\alpha(Q)\overline{F_\alpha(Q)}H(u,u)^d.
\end{equation}
All the endomorphisms and their normalized rational functions
$F_\alpha$ are defined over a single number field: choose
a field containing the coefficients of the finitely many
generators of the CM order, and close it under conjugation.
The uniqueness of the Laurent normalization in C.7.6 gives
the same field of definition for every $F_\alpha$.
Adjoining the hypothesized algebraic number $H(u,u)$ gives
one field for all the values in \eqref{eq:cc-fixedfield}.
The value for $\alpha=0$ is zero, and period shifts do not
change it. Squaring was included because algebraicity of the
individual positive square roots alone would not establish
their membership in one fixed number field.

Logarithmically differentiating the squared CM law at $u$
gives
\[
 a_\alpha=\frac{1}{2\alpha}
     \frac{(d/dz)F_\alpha(\exp_Ez)|_{z=u}}{F_\alpha(Q)}
                  \in\overline{\Q}.
\]
The denominator is nonzero since $Q$ is nontorsion. The
derivation $d/dz$ is the algebraic invariant derivation on
$E$, so the displayed derivative has algebraic value.
Since $D(u)$ is real, conjugation gives the corresponding
formula at $\bar\alpha u$. On a neighborhood where
$\mathcal F\ne0$, its logarithmic gradient and Hessian are
\[
 (\partial_z\log\mathcal F,\partial_w\log\mathcal F)
       =(D(z)-kw,D(w)-kz),\qquad
 \operatorname{Hess}\log\mathcal F=
 \begin{pmatrix}-\wp(z)-s_2&-k\\-k&-\wp(w)-s_2\end{pmatrix}.
\]
Evaluation on $(\alpha,-\bar\alpha)$ proves both identities
in the statement. In particular the Gaussian cancels in the
first derivative and reappears in the second derivative.

To prove $k$ transcendental, recall the real period formula
$\eta_1=(s_2+k)\omega_1$. If $k$ were algebraic, the
nonzero period vector $(\omega_1,\eta_1)$ of the algebraic
universal vector extension of $E$ would lie on the algebraic
line of slope $s_2+k$. Dividing this period by a sufficiently
large integer supplies a nonzero algebraic torsion exponential.
The analytic subgroup theorem
\citep[Theorem 6.1]{huber2022} would give a one-dimensional
algebraic subgroup with that Lie line. Its projection to $E$
is nonzero, so it is an elliptic curve isogenous to $E$ and
splits the vector extension after that isogeny. This is
impossible: its nonzero class in $H^1(E,\mathcal O_E)$ stays
nonzero under an isogeny, whose pullback has a left inverse
after dividing the trace by the degree. These are the same
algebraic logarithm coordinates and extension as in C.7.7;
a fixed algebraic linear change of coordinates leaves the
argument unchanged.

Write the two logarithmic derivatives as
$V_\alpha\log\mathcal F=B_\alpha$ and
$V_\alpha^2\log\mathcal F=A_\alpha+2dk$, where
$A_\alpha,B_\alpha$ are algebraic. Since $H=\mathcal F^2$,
\[
 \frac{V_\alpha^2H(p_\alpha)}{H(p_\alpha)}
       =4B_\alpha^2+2A_\alpha+4dk.
\]
The right side is transcendental, and the asserted conclusion
follows when $H(p_\alpha)\ne0$ is algebraic.
At $\alpha=1$, $a_1=0$, $\wp(u)=4$ and $s_2=-2$,
which give the last displayed formula in the statement.
\end{proof}
